\documentclass[10pt,a4paper]{amsart}
\usepackage[margin=19mm]{geometry}
\usepackage{amsmath,amssymb,mathtools}
\usepackage[scr=rsfs,cal=euler]{mathalfa}
\usepackage{xfrac}
\usepackage[unicode,hidelinks,bookmarks=false]{hyperref}
\newcommand{\ie}{i.e.\ }
\newcommand{\cf}{cf.\ }
\newcommand{\resp}{respectively\ }
\newcommand{\Subsection}{\subsection}
\providecommand{\nobreakdash}{\nobreak\hbox{-}}
\setlength{\emergencystretch}{2em}
\multlinegap0pt
\usepackage{amssymb}
\usepackage{mathtools}
\usepackage[shortlabels]{enumitem}
\newtheorem{theorem}{Theorem}
\newtheorem{proposition}{Proposition}
\newcommand{\RepoTag}{v2.2.8}
\newcommand{\codefile}[1]{%
  \href{https://github.com/SamPetkov/wow284/blob/\RepoTag/#1}{\path{#1}}}
\newcommand{\one}{\mathbf 1}
\newcommand{\resultbox}[1]{\boxed{#1}}
\newcommand{\tr}{\operatorname{tr}}
\title{Counterexamples, Spectral Obstructions, and Deletion Stability for WOW-284}
\author{Samuil Petkov}
\date{Source: arXiv:2607.27452v1 (CC BY 4.0)}
\begin{document}
\maketitle

\noindent\emph{Source and licence.} Counterexamples, Spectral Obstructions, and Deletion Stability for WOW-284. Version arXiv:2607.27452v1 (CC BY 4.0), distributed by arXiv under the Creative Commons Attribution 4.0 International licence, http://creativecommons.org/licenses/by/4.0/, as recorded in the arXiv licence metadata for this exact version and shown on the abstract page https://arxiv.org/abs/2607.27452v1. Full licence notice: \texttt{licenses/arxiv-2607.27452-CC-BY-4.0.txt}.\par\smallskip

\noindent\textit{Editorial scope.} Counted unit: the article's theorem thm:nonadjacent-puncture (source0.tex lines 2237--2263). The article's complete original proof of it is delivered with it (source0.tex lines 2265--2419, one \texttt{proof} environment of 155 lines). No mathematics is written, repaired or re-derived: the statement and the proof are the article's own lines, in the article's own notation, and no retained line is substituted or re-flowed.  Retained uncounted as the source-local context the proof needs: lines 2421--2440; lines 2457--2472.  Nothing was omitted from any retained span, so the package carries no omission record.  No retained line contains a citation, so the wrapper carries no bibliography and no bibliography entry.  No retained line contains a figure environment, an image-inclusion command or drawing code, so no image is delivered and the package declares no asset dependency.  Editorial formatting only, in addition: an AMS \texttt{amsart} preamble that restates the article's own declarations and macros used by the retained lines, namely the environments \texttt{theorem}, \texttt{proposition} on separate counters, together with the standard packages those lines need.  The retained environments print in this document as Theorem 1, Proposition 1 in order of appearance, whereas the article prints the same items with the numbering of its own full document.  The complete native source of the article as distributed in the arXiv e-print for this exact version is bundled unchanged as \texttt{sources/206349/source0.tex}.
\begin{proposition}[Deletion stability]\label{prop:deletion-stability}
Let \(G\) be connected with at least two vertices, let
\(S\subseteq V(G)\), and suppose \(H=G-S\) is connected with at least two
vertices.  Put
\[
 D_0=D(G)[V(H)],
 \qquad E_S=D(H)-D_0,
\]
\[
 a=\delta^*(G),
 \quad b=\delta^*(H),
 \quad \gamma=\Phi(G).
\]
Then
\[
 \resultbox{
 \Phi(H)\ge\gamma-(a-b)+\lambda_{\min}(E_S).
 }
\]
\end{proposition}

\begin{theorem}[Small-puncture normal form]\label{thm:small-puncture}
Let \(M\) be a degree-\(k\) Moore graph of diameter two, let
\(S\subseteq V(M)\) have size \(s\le k-1\), and put \(H=M-S\).  Then \(H\)
is connected, has diameter at most three, and
\[
 \resultbox{\delta^*(H)=k-\frac{s}{k}.}
\]
Let \(B\) be the surviving-vertex by deleted-vertex incidence matrix, with
\(B_{xz}=1\) exactly when \(x\sim_M z\).  Then
\[
 \resultbox{
 D(H)=2(J-I)-A(H)+BB^{\mathsf T}
 -\operatorname{diag}(BB^{\mathsf T}).
 }
\]
\end{theorem}

\begin{theorem}[Two nonadjacent deleted vertices]\label{thm:nonadjacent-puncture}
Let \(u,v\) be nonadjacent vertices of \(M\), \(k\ge5\), and put
\(H=M-\{u,v\}\).  Define
\begin{align*}
R_k(x)={}&x^4+(10-2k^2)x^3+(2k^3-17k^2-2k+36)x^2\\
&+(12k^3-49k^2-4k+53)x\\
&-2k^4+17k^3-38k^2+5k+20,
\end{align*}
\[
 M_- =\frac{k(k-2)+(k^2-4k+2)\Delta}{2\Delta},
 \quad
 M_+ =\frac{-k(k-2)+(k^2-4k+2)\Delta}{2\Delta}.
\]
Then
\begin{align*}
\chi_{D(H)}(x)={}&(x-k+3)R_k(x)
 (x^2+4x-k+3)^{k-2}\\
&\cdot(x^2+4x-k+5)^{k-2}
 \left(x+\frac{\Delta+3}{2}\right)^{M_-}\\
&\cdot\left(x-\frac{\Delta-3}{2}\right)^{M_+}.
\end{align*}
Moreover,
\[
 \delta^*(H)=k-\frac2k,
\]
and \(H\) is a strict counterexample for every realizable integer \(k\ge6\).
\end{theorem}

\begin{proof}
The deleted vertices have a unique common neighbour \(w\).  With
\(A=N(u)\setminus\{w\}\), \(B=N(v)\setminus\{w\}\),
\(C=N(w)\setminus\{u,v\}\), and
\[
 Z=V(M)\setminus\bigl(\{u,v,w\}\cup A\cup B\cup C\bigr),
\]
the surviving graph has cell sizes, in the order \(\{w\},A,B,C,Z\),
\[
 1,\quad k-1,\quad k-1,\quad k-2,\quad (k-1)(k-2).
\]
The Moore common-neighbour rule makes the \(A\)--\(B\) edges a perfect
matching and the five-cell partition equitable.  It also supplies a
length-three replacement path for every pair whose unique length-two path
used \(u\) or \(v\); hence every such new distance is exactly three.  On the
cell-constant space, the row-sum distance quotient is
\[
\begin{pmatrix}
0&3k-3&3k-3&k-2&2(k-1)(k-2)\\
3&3k-6&2k-3&2k-4&2k^2-7k+6\\
3&2k-3&3k-6&2k-4&2k^2-7k+6\\
1&2k-2&2k-2&2k-6&2k^2-7k+5\\
2&2k-3&2k-3&2k-5&2k^2-7k+5
\end{pmatrix}.
\]
Its characteristic polynomial is \((x-k+3)R_k(x)\).

Identify \(\mathbb R^A\) with \(\mathbb R^B\) through the perfect matching,
and let \(R_A,R_B,R_C\) be the incidence matrices from the three
nonconstant cells to \(Z\).  Let \(T\) be the adjacency matrix on \(Z\).
The Moore identity gives
\[
\begin{gathered}
R_AR_A^{\mathsf T}=R_BR_B^{\mathsf T}=(k-2)I,\quad
R_CR_C^{\mathsf T}=(k-1)I,\\
R_AR_B^{\mathsf T}=J-I,\quad
R_AR_C^{\mathsf T}=R_BR_C^{\mathsf T}=J,\\
R_AT+R_B=J-R_A,\quad
R_BT+R_A=J-R_B,\quad
R_CT=J-R_C,\\
R_A^{\mathsf T}R_A+R_B^{\mathsf T}R_B+R_C^{\mathsf T}R_C+T^2
=(k-1)I-T+J.
\end{gathered}
\]
Moreover, every vertex of \(Z\) has one neighbour in each of \(A,B,C\),
so
\[
 R_A^{\mathsf T}\one_A
 =R_B^{\mathsf T}\one_B
 =R_C^{\mathsf T}\one_C
 =\one_Z,
 \qquad T\one_Z=(k-3)\one_Z.
\]

Use the coordinate order \((\{w\},A,B,C,Z)\).  Choose orthogonal bases
\(\mathcal B_A\) of \(\one_A^\perp\) and
\(\mathcal B_C\) of \(\one_C^\perp\).  For
\(\xi\in\mathcal B_A\) and \(\eta\in\mathcal B_C\), put
\[
\begin{aligned}
u_\xi^\pm&=(0,\xi,\pm\xi,0,0),&
v_\xi^\pm&=(0,0,0,0,
 (R_A^{\mathsf T}\pm R_B^{\mathsf T})\xi),\\
u_\eta^C&=(0,0,0,\eta,0),&
v_\eta^C&=(0,0,0,0,R_C^{\mathsf T}\eta),
\end{aligned}
\]
and define
\[
 W_\xi^\pm=\operatorname{span}\{u_\xi^\pm,v_\xi^\pm\},
 \qquad
 W_\eta^C=\operatorname{span}\{u_\eta^C,v_\eta^C\}.
\]
The displayed identities show that the relevant incidence maps are
injective, that these two-dimensional modules are mutually orthogonal,
and that they are orthogonal to the five-dimensional cell-constant space
\(W_{\mathrm{const}}\).  For
\[
 K=\ker R_A\cap\ker R_B\cap\ker R_C
\]
one has
\[
 K\subseteq\one_Z^\perp,\qquad
 \dim K=(k-2)(k-4),\qquad T(K)\subseteq K.
\]
Thus \(W_{\mathrm{const}}\), the modules \(W_\xi^\pm,W_\eta^C\), and
\(0\oplus0\oplus0\oplus0\oplus K\) form an orthogonal direct sum, since
\[
 5+4(k-2)+2(k-3)+(k-2)(k-4)=k^2-1.
\]

In the ordered generator pairs
\((u_\xi^+,v_\xi^+)\), \((u_\xi^-,v_\xi^-)\), and
\((u_\eta^C,v_\eta^C)\), respectively, the distance operator has matrices
\[
 \begin{pmatrix}-4&-(k-3)\\-1&0\end{pmatrix},\qquad
 \begin{pmatrix}-2&-(k-1)\\-1&-2\end{pmatrix},\qquad
 \begin{pmatrix}-2&-(k-1)\\-1&-1\end{pmatrix}.
\]
These give the two displayed quadratic factors and \(k-3\) copies of
each Moore linear factor.  On \(K\), the \(Z\)-block identity gives
\[
 T^2+T-(k-1)I=0,
 \qquad D(H)|_K=-2I-T.
\]
The constant direction of \(Z\) has \(T\)-eigenvalue \(k-3\), while the
\(Z\)-images in the symmetric, antisymmetric, and common-neighbour
modules have eigenvalues \(-2,0,-1\), with dimensions
\(k-2,k-2,k-3\), respectively.  Since \(\tr T=0\),
\[
 \tr(T|_K)=2(k-2).
\]
Dimension and trace give the residual multiplicities; after adding the
common-neighbour copies, they are precisely \(M_-\) and \(M_+\).

The value \(\delta^*(H)=k-\frac2k\) is again the \(s=2\) case of
Theorem~\ref{thm:small-puncture}.  For strictness, put
\[
 E_S=D(H)-D(M)[V(H)].
\]
Up to a simultaneous permutation of rows and columns, \(E_S\) is the
adjacency matrix of two copies of \(K_k\) meeting in \(w\), together with
\(k(k-2)\) isolated vertices.  Hence
\[
 \lambda_{\min}(E_S)=\frac{k-2-\sqrt{k^2+4k-4}}2.
\]
Since the parent Moore graph has score
\(k-(3+\Delta)/2\) and deletion lowers the minimum dual degree by
\(2/k\),
Proposition~\ref{prop:deletion-stability} gives
\[
 \Phi(H)\ge
 \frac{3k-5-\Delta-\sqrt{k^2+4k-4}}2-\frac2k.
\]
For \(k\ge6\),
\[
\Delta<2\sqrt{k}-\frac{3}{4\sqrt{k}},
\qquad
\sqrt{k^2+4k-4}<k+2-\frac4{k+2}.
\]
The lower bound is therefore greater than
\[
 f(k)=k-\frac72-\sqrt{k}+\frac{3}{8\sqrt{k}}
       +\frac2{k+2}-\frac2k.
\]
Now \(f(6)=29/12-15\sqrt6/16>0\), since
\(116^2>6\cdot45^2\).  Moreover \(f'(x)>0\) for \(x\ge6\):
the absolute values of the three negative terms in \(f'(x)\) have sum less than
\(1/4+1/64+1/32<1\), while \(2/x^2>0\).  Thus \(\Phi(H)>0\).
The direct-sum decomposition, injectivity, orthogonality, all recomputed
distances, and the strict comparison
are independently checked by
\codefile{scripts/verify_proof_audit_03_nonadjacent_puncture.py} and
\codefile{scripts/verify_research_extensions_exact.py}.
\end{proof}

\end{document}
